\chapter{100 câu cục súc về tình yêu, rung động và tan vỡ tuổi trẻ}
\markboth{100 câu cục súc về tình yêu, rung động và tan vỡ tuổi trẻ}{100 câu cục súc về tình yêu, rung động và tan vỡ tuổi trẻ}

\section{25 câu về rung động và ảo tưởng ban đầu}
\begin{enumerate}[leftmargin=*]
\item Thích một người không làm em sâu hơn tự động.
\item Nhiều rung động đầu đời đẹp vì nó chưa đụng đời đủ mạnh.
\item Đừng gọi là định mệnh chỉ vì hai đứa nhắn tin hợp vài đêm.
\item Cảm giác hợp chưa chắc đi được đường dài.
\item Người làm tim em đập nhanh chưa chắc là người làm đời em yên.
\item Tuổi trẻ rất giỏi nhầm giữa được chú ý và được yêu.
\item Có những mối bắt đầu vì cô đơn chứ không vì hiểu nhau.
\item Đẹp đôi ngoài ảnh không đồng nghĩa đỡ làm nhau mệt ở ngoài đời.
\item Nhiều đứa yêu người ta thật ra là yêu cảm giác mình quan trọng trong mắt người ta.
\item Đừng vì một ánh mắt mà viết cả tương lai.
\item Cái gì đến quá nhanh thường cần nhìn chậm lại.
\item Tình cảm thật không sợ rõ ràng.
\item Nếu một mối quan hệ chỉ sống bằng đoán ý thì sớm muộn cũng ngộp.
\item Yêu không làm hết nỗi bất an cũ nếu em chưa chịu chữa nó.
\item Nhiều người tìm người yêu để né việc học ở một mình.
\item Tim rung không phải lúc nào não cũng nên tắt.
\item Đừng vì được một người chọn mà quên tự chọn mình.
\item Ai khiến em phải diễn quá nhiều để được thích thì nên xem lại.
\item Rung động đẹp nhất là khi nó không làm em đánh rơi phẩm giá.
\item Đừng gọi mọi thứ mãnh liệt là sâu.
\item Có những mối hợp vui nhưng không hợp lớn lên.
\item Yêu ai đó không phải giấy phép để bỏ bê cả đời mình.
\item Người xứng đáng với em sẽ không cần em tự nhỏ lại quá nhiều.
\item Thích là chuyện tim, ở lại hay không còn là chuyện não.
\item Tỉnh táo trong yêu không làm mất lãng mạn, nó giữ em khỏi ngu.
\end{enumerate}

\section{25 câu về yêu mà mất mình}
\begin{enumerate}[leftmargin=*,resume]
\item Yêu tới mức bỏ hết bạn bè, việc học, giấc ngủ thì nghe không ổn lắm.
\item Một mối quan hệ bắt em đánh đổi nhân phẩm không đáng để giữ.
\item Đừng gọi là yêu sâu nếu em đang sống bằng lo âu và dò đoán.
\item Có người không yêu em, họ chỉ thích cảm giác được em xoay quanh.
\item Mất mình trong tình cảm không phải chung thủy, là thiếu trục sống.
\item Đừng xem kiểm soát là quan tâm phiên bản mạnh tay.
\item Bị giữ bằng sợ hãi không phải được yêu.
\item Nếu lúc nào em cũng phải chứng minh giá trị để được ở lại thì em đang ở sai chỗ.
\item Yêu mà cứ phải cầu xin sự tử tế thì nghỉ đi cho đỡ nhục tim.
\item Nhiều người chấp nhận bị đối xử tệ chỉ vì sợ quay lại cô đơn.
\item Đừng đổi bình yên lấy vài tin nhắn ngọt đúng lúc.
\item Yêu không nên làm em ghét bản thân hơn.
\item Một tình cảm khiến em phải luôn thắc thỏm xem mình có đủ không thường rất bào mòn.
\item Có những người nói yêu như cách giữ một món đồ khỏi tay người khác.
\item Tự hạ giá mình để giữ người khác là khoản lỗ rất ngu.
\item Đừng dùng chịu đựng để chứng minh tình yêu.
\item Người trưởng thành biết yêu mà không nuốt luôn ranh giới của mình.
\item Nếu một mối làm em kiệt quệ lâu ngày, nó không còn đáng thơ đâu.
\item Yêu ai đó không phải bỏ quyền được tôn trọng.
\item Càng quên mình trong tình cảm càng dễ oán về sau.
\item Muốn thương lâu phải còn giữ được phần nào tỉnh táo.
\item Đừng lấy câu ai yêu nhiều hơn người đó thua để hợp thức hóa việc mình bị đối xử tệ.
\item Có những thứ em tưởng là hy sinh, thật ra là tự bỏ rơi.
\item Mối quan hệ lành mạnh không làm em sống như đang thi nói dối và đoán lòng suốt ngày.
\item Yêu mà không còn dám là mình thì đang lạc rồi.
\end{enumerate}

\section{25 câu về chia tay và hậu đau lòng}
\begin{enumerate}[leftmargin=*,resume]
\item Chia tay không giết em, nhưng nó lột ra nhiều chỗ em từng bám quá chặt.
\item Đau không đáng ngại bằng việc em lấy nỗi đau làm căn cước mới.
\item Mất một người không có nghĩa là mất luôn khả năng được yêu.
\item Có những cuộc chia tay là phẫu thuật chứ không phải tai nạn.
\item Đừng níu một người chỉ vì em chưa quen khoảng trống.
\item Cái đau sau chia tay không chỉ vì mất người, còn vì mất thói quen.
\item Nỗi nhớ rất giỏi tô hồng những ngày cũ vốn không hề dễ chịu vậy.
\item Đừng gọi mọi kết thúc là thất bại, nhiều cái kết là cứu mạng.
\item Có người buồn vì yêu sâu, có người buồn vì cái tôi bị bỏ.
\item Người cũ không phải cái cớ để em hỏng thêm.
\item Chia tay xong đừng biến đời mình thành hiện trường cần người kia quay lại mới sống được.
\item Buồn thì buồn cho tới nơi, nhưng đừng tự làm nhục mình trong lúc buồn.
\item Có những tin nhắn không gửi đi là một dạng trưởng thành.
\item Đừng xin lại thứ người ta đã thả tay quá rõ.
\item Nếu một người rời đi mà em tan quá mạnh, có thể em từng đặt hết giá trị mình vào họ.
\item Hậu chia tay dạy em cách cai nghiện cảm xúc rất đau nhưng rất thật.
\item Không phải ai quay lại cũng vì yêu, có người chỉ cô đơn đúng lúc.
\item Tự trọng sau chia tay quý hơn mấy pha níu kéo lên mạng rất nhiều.
\item Nước mắt không làm em yếu, nhưng tự hủy sau nước mắt thì có.
\item Có những người em phải thôi nhớ theo kiểu muốn quay về, và bắt đầu nhớ theo kiểu biết ơn bài học.
\item Đừng làm người mới trả tiền cho vết thương người cũ.
\item Hồi phục không thẳng tắp, có ngày ổn có ngày dở, miễn là đừng quay lại chỗ cũ chỉ vì nhớ.
\item Người bỏ em có thể từng là mảnh đẹp, nhưng không còn là mái nhà đúng cho em.
\item Cái đau rồi qua, chỉ cần em đừng tự thêu lại nó mỗi đêm.
\item Một trái tim đã từng vỡ vẫn có thể học cách yêu tử tế hơn chứ không nhất thiết chai đi.
\end{enumerate}

\section{25 câu về yêu tử tế và lớn lên cùng nhau}
\begin{enumerate}[leftmargin=*,resume]
\item Yêu đẹp là khi cả hai đều lớn hơn, không phải một người mòn đi cho người kia sáng.
\item Tình cảm bền không cần ồn, chỉ cần đủ thật và đủ rõ.
\item Người yêu tử tế không bắt em đoán quá nhiều.
\item Quan tâm đúng là làm cuộc sống nhau đỡ nặng, không làm nhau nặng thêm.
\item Tình yêu tốt không phá nhịp học, nhịp sống và lòng tự trọng của em.
\item Một mối quan hệ đáng giữ sẽ chịu nổi sự rõ ràng.
\item Yêu ai đó mà vẫn tôn trọng thời gian và mục tiêu của nhau mới là trưởng thành.
\item Chín chắn trong yêu là biết dỗ, biết nói, biết dừng khi nóng.
\item Người hợp với em không nhất thiết giống em, nhưng họ sẽ không xem phần cốt lõi của em như phiền phức.
\item Tình cảm bền thường xây bằng thói quen tử tế nhỏ, không chỉ bằng cảm xúc lớn.
\item Tôn trọng nhau lúc giận mới là nơi tình yêu lộ chất.
\item Được yêu đúng cách làm em yên chứ không làm em ám ảnh.
\item Một người tốt cho em thường không khiến em phải hy sinh phẩm giá để ở lại.
\item Yêu là cùng nhìn về hướng sống, không chỉ nhìn nhau mờ mắt.
\item Nói chuyện được với nhau lúc khó còn quý hơn rất nhiều lời ngọt lúc vui.
\item Tình cảm sâu không ngại xin lỗi, cũng không ngại sửa.
\item Yêu mà vẫn giữ được bạn bè, việc học và ước mơ là một kiểu cân bằng đẹp.
\item Người trưởng thành biết khi nào cần ôm, khi nào cần để nhau yên.
\item Hợp nhau một phần là may, phần còn lại là nỗ lực biết điều.
\item Đừng tìm người hoàn hảo, tìm người đủ thật để cùng sửa những phần chưa ổn.
\item Tình yêu không chữa hết đời em, nhưng có thể làm đoạn đường đỡ cô độc nếu đặt đúng người.
\item Một mối quan hệ lành mạnh không cần em biến thành người khác để tồn tại.
\item Người yêu tốt giúp em sống tử tế hơn với chính mình.
\item Yêu đúng người không làm đời bớt khó hoàn toàn, nhưng làm em bớt phải thủ sẵn áo giáp ở mọi phút.
\item Điều đẹp nhất của tình yêu là hai người cùng đỡ tệ hơn, chứ không cùng dở đi.
\end{enumerate}

\chapter{100 câu cục súc về kỷ luật, bản lĩnh và đứng thẳng giữa đời}
\markboth{100 câu cục súc về kỷ luật, bản lĩnh và đứng thẳng giữa đời}{100 câu cục súc về kỷ luật, bản lĩnh và đứng thẳng giữa đời}

\section{25 câu về tự quản và nề nếp}
\begin{enumerate}[leftmargin=*]
\item Không ai nhìn mà em vẫn làm đúng là bản lĩnh thật.
\item Kỷ luật nghe khô, nhưng người thiếu nó thường sống rất nhão.
\item Nề nếp không bóp chết tự do, nó giữ tự do khỏi biến thành bừa bộn.
\item Đúng giờ là tôn trọng người khác và bớt khinh chính mình.
\item Người không quản nổi giờ ngủ rất khó quản được những giấc mơ lớn.
\item Cái gì làm đều được mới đáng gọi là năng lực.
\item Đừng gọi mình phóng khoáng nếu em chỉ đang sống thiếu trật tự.
\item Có kỷ luật rồi em mới thấy trước đây mình thất thoát bao nhiêu sức vô lý.
\item Bản lĩnh thường nhìn chán vì nó sống trong việc lặp lại.
\item Cứ làm điều cần làm trước rồi cảm xúc tính sau.
\item Người đáng tin không được tạo bởi vài lần cao hứng.
\item Muốn đời bớt loạn thì bàn học, giấc ngủ, lịch sống cũng phải bớt loạn.
\item Tự quản kém làm mọi mục tiêu đều thành chuyện may rủi.
\item Kỷ luật không làm em ngầu ngay, nó làm em bớt yếu trước đã.
\item Đừng để tâm trạng làm sếp quá lâu.
\item Làm việc đúng lúc là cách giảm rất nhiều nỗi sợ vô hình.
\item Người bền ít khi đợi hứng, họ đợi thói quen tới giờ.
\item Từ chối vài cái vui ngắn là cách mở đường cho vài cái đáng hơn.
\item Kỷ luật là nói không với phiên bản dễ dãi của mình mỗi ngày.
\item Sống có hệ thống thì lúc gặp biến mới đỡ tan.
\item Không ai mạnh trong một ngày, người ta mạnh trong cách sống lặp lại.
\item Sắp xếp được mình là đã thắng được kha khá hỗn loạn ngoài kia.
\item Kỷ luật không sexy trên mạng, nhưng rất có giá ngoài đời.
\item Những người đi xa hiếm khi sống kiểu tiện đâu theo đó.
\item Nề nếp là một dạng yêu tương lai của mình bằng hành động.
\end{enumerate}

\section{25 câu về bản lĩnh khi bị thử lửa}
\begin{enumerate}[leftmargin=*,resume]
\item Bản lĩnh không lộ ra lúc em được chiều, nó lộ lúc em bị ép góc.
\item Căng quá mới biết mình luyện được bao nhiêu.
\item Đừng nhận mạnh nếu chỉ cần bị chê chút đã gãy.
\item Người vững không phải người không run, mà là người run vẫn không bỏ vị trí.
\item Bản lĩnh là biết đau nhưng không bẻ cong nguyên tắc quá nhanh.
\item Cái gì em giữ được trong lúc khó mới là cái thật sự thuộc về em.
\item Đời đẩy mạnh một cái là nhiều mặt nạ rơi rất nhanh.
\item Bị hiểu lầm mà không tự phát nổ ngay là một loại mạnh hiếm.
\item Người bản lĩnh không cần thắng khẩu chiến mọi lúc.
\item Có khi im đúng lúc khó hơn cãi cho sướng miệng.
\item Thua mà không biến thành người cay đắng là trình độ.
\item Khi mọi thứ không theo ý mình mà vẫn giữ được tử tế mới là chất.
\item Dám chịu hậu quả cho quyết định của mình là bản lĩnh.
\item Người mạnh không dựa hoàn toàn vào việc luôn được công nhận.
\item Càng ít bản lĩnh càng thích phô giọng dữ.
\item Bản lĩnh và cứng đầu nhìn gần giống, khác ở chỗ một bên biết sửa khi thấy sai.
\item Muốn mạnh phải chịu được vài lần bị đời bóc trần.
\item Không phải cú ngã nào cũng nên kể than, có cú để tự mình tiêu hóa sẽ lớn hơn.
\item Bản lĩnh là dám giữ cái đúng ngay cả khi nó làm em lạc quẻ.
\item Người có lực thường không cần bày vẻ bất cần quá nhiều.
\item Giữ lời hứa với người khác quý, giữ lời hứa với mình còn khó hơn.
\item Cái khó không phải mạnh một lúc, mà là không thối giữa chừng.
\item Đừng luyện vẻ lì, hãy luyện nội lực.
\item Trưởng thành hơn là khi em không còn phải làm ồn để chứng minh mình không sợ.
\item Một người đứng thẳng được sau vài lần bẹp là người đáng nể thật.
\end{enumerate}

\section{25 câu về nguyên tắc và ranh giới}
\begin{enumerate}[leftmargin=*,resume]
\item Có nguyên tắc không làm em khó gần, nó chỉ làm em khó bị dùng.
\item Ranh giới là hàng rào bảo vệ phần tử tế của em khỏi bị bào mòn.
\item Ai cũng muốn em dễ chịu, ít ai vui khi em rõ ràng.
\item Người không tôn trọng ranh giới của em thường là người hưởng lợi khi em không có ranh giới.
\item Nguyên tắc chỉ có giá khi nó tồn tại lúc em bị cám dỗ.
\item Sống dễ quá với ai cũng thường khó với chính mình về sau.
\item Đừng để lòng thương làm em mềm đến mức mất xương.
\item Tử tế không phải giấy phép cho người khác đi giày lên mặt mình.
\item Một câu không nói đúng lúc có thể cứu nhiều tháng bực âm thầm.
\item Rõ ràng sớm thường đỡ đau hơn nhẫn nhịn lâu rồi bùng nổ.
\item Có những việc em không làm không phải vì không dám, mà vì không đáng.
\item Người lớn thực sự biết phân biệt giữa linh hoạt và đánh mất cốt lõi.
\item Càng quý mình càng ít thích mặc cả nhân phẩm.
\item Không phải lời mời nào cũng đáng gật.
\item Cần được yêu thích quá mức khiến em rất dễ bán nguyên tắc.
\item Nói không tử tế là kỹ năng sống, không phải thô lỗ.
\item Ranh giới rõ giúp người tử tế ở lại và người lợi dụng tự lộ.
\item Đừng đợi bị dồn tới nghẹt mới nhớ mình có quyền từ chối.
\item Không thuận hết lòng người không phải thất bại.
\item Nguyên tắc tốt giúp em đỡ phải quyết định lại từ đầu mỗi lần yếu lòng.
\item Có xương sống rồi em mới thương được lâu mà không oán.
\item Rõ ràng không giết mối quan hệ đúng, nó chỉ giết mối quan hệ mập mờ sống nhờ sự dễ dãi của em.
\item Ai khó chịu vì em lành mạnh hơn có lẽ từng quen hưởng lợi từ phần chưa lành của em.
\item Đừng xem ranh giới là bức tường, hãy xem nó là cánh cổng có chọn lọc.
\item Người biết giữ mình sẽ đỡ phải đi nhặt mình sau này.
\end{enumerate}

\section{25 câu về đứng thẳng giữa đời xô lệch}
\begin{enumerate}[leftmargin=*,resume]
\item Đời lệch nhiều chỗ, nhưng em không cần cong theo hết.
\item Không phải vì ai cũng lách mà mình phải xem lách là khôn.
\item Sống đàng hoàng đôi khi thiệt trước, nhưng thường đỡ nát sau.
\item Cái khó là giữ chuẩn trong môi trường thấp chuẩn.
\item Đừng lấy sự phổ biến của cái sai để làm nó hợp lý.
\item Có những lúc đứng một mình vẫn đáng hơn đứng sai với số đông.
\item Thấy người khác ăn gian mà mình vẫn giữ sạch là một loại mạnh không ồn.
\item Đừng trách mình ngây khi chọn tử tế, chỉ cần đừng tử tế thiếu não.
\item Bị coi là cứng có khi chỉ vì em không chịu trôi theo thói dễ dãi chung.
\item Người biết giữ đường thường đi chậm hơn đám chen lấn lúc đầu.
\item Nếu không có lõi, mỗi cơn gió dư luận đều có thể bẻ em.
\item Đứng thẳng không hứa đời dễ, nhưng giúp em ngủ yên hơn.
\item Giữ sạch giữa chỗ bẩn không phải để làm thánh, mà để không ghét mình.
\item Ai cũng có giá của mình, đừng để nó rẻ đến mức chính em ngại nhìn.
\item Càng sống lâu càng thấy bình yên với lương tâm là một loại sang hiếm.
\item Đúng không luôn dễ nhận ra, nhưng sai nhiều khi rất dễ nếu em chịu ngồi yên vài phút.
\item Có những chiến thắng nhìn ngoài rất nhỏ nhưng bên trong cứu được cả tự trọng.
\item Đừng chê người tử tế là khờ khi em còn chưa thử sống đủ sạch để hiểu cái giá và giá trị của nó.
\item Đời xô đẩy em dữ, nhưng cái hướng em chọn vẫn nói lên em là ai.
\item Thẳng lưng đi chậm còn hơn cúi đầu đi nhanh mãi.
\item Người trẻ đáng nể không phải người có vẻ bất cần, mà là người có chuẩn riêng và chịu giữ.
\item Bản lĩnh cuối cùng không nằm ở thắng ai, mà ở việc không phản bội mình trước áp lực.
\item Dư luận ồn hôm nay mai sẽ quên, lương tâm em thì không quên nhanh vậy đâu.
\item Có nội lực rồi em không cần mượn đám đông để cảm thấy mình đúng.
\item Đời có thể không thưởng ngay cho sự đàng hoàng, nhưng nó thường tính sổ rất rõ với sự cẩu thả nhân cách.
\end{enumerate}

\chapter{100 câu cục súc về vấp ngã, thất bại và làm lại từ đầu}
\markboth{100 câu cục súc về vấp ngã, thất bại và làm lại từ đầu}{100 câu cục súc về vấp ngã, thất bại và làm lại từ đầu}

\section{25 câu về cú ngã đầu đời}
\begin{enumerate}[leftmargin=*]
\item Ngã lần đầu đau vì em tưởng mình khác với việc phải ngã.
\item Thất bại không hỏi em đã sẵn sàng chưa.
\item Đời vả một cái thường không để làm nhục em, mà để chỉnh góc nhìn.
\item Có những lần thua chỉ để giết bớt ảo tưởng rất dày.
\item Đừng xem cú ngã là bằng chứng em vô dụng.
\item Người chưa từng đau thường đánh giá đời rất non.
\item Mới thua một trận mà kết luận cả đời thì hơi vội.
\item Cú ngã nào cũng có hai phần: cái mất và cái lộ ra.
\item Nhiều khi em đau không chỉ vì thất bại, mà vì hình ảnh mình tự dựng bị bể.
\item Thua sớm một chút đôi khi cứu em khỏi tự mãn lâu.
\item Đừng vội ghét đời khi đời chỉ đang nói em cần học thêm.
\item Có những cú ngã không đẹp nhưng rất đúng thời điểm.
\item Vấp không đáng shame bằng cố đi tiếp với cái sai y nguyên.
\item Càng sợ ngã càng khó đi đường lớn.
\item Muốn lớn mà không muốn đau là đòi hỏi hơi ngây thơ.
\item Thất bại đầu đời thường lộn xộn và hơi quê, nhưng nó dạy nhanh.
\item Chưa ai học được bản lĩnh chỉ bằng lời động viên ngọt.
\item Đừng đòi cuộc sống nhẹ tay chỉ vì em đang mỏng.
\item Những cú vấp thật đôi khi giúp em thôi sống giả.
\item Không ai trưởng thành chỉ từ những lần được chiều.
\item Cú ngã không nhỏ hay lớn bằng cách em tiêu hóa nó.
\item Một người trẻ biết nhìn thẳng vào cú thua đã hơn rất nhiều người chỉ biết kể khổ.
\item Đau là học phí, đừng biến nó thành bản sắc.
\item Có thể buồn sâu, nhưng đừng tự xem mình hết đường.
\item Ngã không định nghĩa em, cách em đứng dậy mới nói nhiều hơn.
\end{enumerate}

\section{25 câu về đối diện thất bại cho ra người}
\begin{enumerate}[leftmargin=*,resume]
\item Thất bại không cần em đẹp, nó cần em thật.
\item Đừng kể lể hay quá nếu cuối cùng vẫn không chịu sửa gì.
\item Chấp nhận thua không phải đầu hàng, mà là dừng cãi với sự thật.
\item Càng sớm nhận phần lỗi của mình càng sớm hết bất lực.
\item Người đổ lỗi giỏi thường hồi phục chậm.
\item Có những cú thua phải ngồi một mình rất lâu mới hiểu nổi.
\item Tự thương hại quá mức là cách kéo dài hậu quả.
\item Thất bại không làm em nhỏ đi, phản ứng trẻ con với thất bại mới làm vậy.
\item Đau xong rồi phải nhìn số liệu, nhìn nguyên nhân, nhìn cách.
\item Đừng cố tìm thủ phạm khi thứ em cần nhất là bài học.
\item Thua mà còn giữ được sự tử tế là tốt, nhưng thua mà còn giữ được khả năng học mới quan trọng hơn.
\item Cứ hỏi tại sao là nạn tôi mãi thì em sẽ bỏ lỡ câu bây giờ làm gì.
\item Đời không hứa công bằng tuyệt đối, nhưng nó vẫn để nhiều cơ hội cho người chịu chỉnh mình.
\item Thất bại là tấm gương xấu xí nhưng rất trung thực.
\item Càng cay cú với kết quả càng nên chậm lại phân tích.
\item Đừng hô làm lại nếu em chưa chịu mổ xác cái cũ.
\item Có những lần thua vì thiếu năng lực, có lần vì thiếu bền, đừng gộp tất cả thành xui.
\item Mạnh không phải khóc ít, mà là khóc xong biết làm gì tiếp.
\item Đừng biến một cú hụt thành lý do cho năm tháng thả trôi.
\item Muốn ngoi lên phải chịu nhìn mình lúc bết nhất.
\item Phủ nhận sự thật chỉ làm bài học đổi tên chứ không biến mất.
\item Ai cũng muốn comeback đẹp, ít ai chịu phần mò mẫm rất xấu trước đó.
\item Thua mà bớt nói nhiều, làm kỹ hơn thì đáng nể.
\item Thất bại không cần em diễn mạnh, nó cần em thành thật với mức độ yếu của mình.
\item Học được từ cú thua là cách duy nhất để nó không phí hoàn toàn.
\end{enumerate}

\section{25 câu về làm lại và đi tiếp}
\begin{enumerate}[leftmargin=*,resume]
\item Làm lại không hào hùng lắm, nó thường lặng và khá nhục.
\item Muốn đi tiếp phải chấp nhận mình đang ở đâu trước đã.
\item Bắt đầu lại ở vị trí thấp hơn không đáng xấu hổ bằng cố giữ cái vỏ cao đã rỗng.
\item Đi tiếp chậm còn hơn đứng đó chửi quá khứ.
\item Làm lại giỏi là kỹ năng sống ít được dạy nhưng ai cũng cần.
\item Cái khó nhất của khởi động lại là nuốt cái tôi.
\item Nếu còn thở còn học được, nghe cliché nhưng vẫn đúng.
\item Đừng đợi đủ tự tin mới làm lại, làm lại rồi tự tin mới mọc dần.
\item Mỗi lần đứng lên tệ tệ vẫn hơn một lần nằm xuống rất thơ.
\item Bắt đầu nhỏ không làm mục tiêu em nhỏ.
\item Đi từng bước xấu xấu nhưng thật còn hơn ngồi mơ màn trở lại hoành tráng.
\item Hồi phục không phải chạy thật nhanh, là không bỏ cuộc với hướng đúng.
\item Có hôm chỉ cần không trượt lại cũng là tiến bộ.
\item Làm lại thành công thường nhìn rất tẻ nếu quay phim không có nhạc nền.
\item Đừng ngại về số 0 nếu số 10 cũ của em là đồ giả.
\item Đi tiếp cần ý chí, nhưng cũng cần hệ thống đỡ ngu hơn trước.
\item Hứa ít lại, quay về làm những thứ nền trước.
\item Ai từng nát mà đứng dậy được thường có chiều sâu khác người chỉ biết thành công dễ.
\item Làm lại không xóa sạch quá khứ, nhưng cho nó vai khác trong đời em.
\item Cứ đều một ít rồi em sẽ thấy mình không còn xa người mình muốn thành như tưởng tượng.
\item Đừng thần tượng cú bứt phá, hãy học sự bền bỉ rất không bắt mắt.
\item Người kiên nhẫn với bản thân sau thất bại đi xa hơn người chỉ biết tự đập mình.
\item Tự tha thứ đúng cách là để làm tiếp tử tế hơn, không phải để thôi cố.
\item Hôm nay bớt dở hơn hôm qua một chút đã là đường đi rồi.
\item Chuyện đẹp nhất của một cú ngã là nó không ngăn được em đi tiếp mãi.
\end{enumerate}

\section{25 câu về biến đau thành chiều sâu}
\begin{enumerate}[leftmargin=*,resume]
\item Không phải cứ đau là sâu, phải tiêu hóa được nó mới sâu.
\item Vết thương nào cũng có thể biến em thành người tinh tế hơn hoặc cay độc hơn.
\item Đau được dùng đúng cách sẽ dạy em nhìn người hiền hơn.
\item Trải qua chuyện xấu mà không thành người xấu hơn là thành tựu lớn.
\item Nỗi đau qua lọc sẽ thành trực giác tốt.
\item Có chiều sâu không phải nói câu nào cũng buồn, mà là hiểu cái giá của nhiều thứ.
\item Sau vài lần gãy em sẽ biết quý cái bình thường hơn.
\item Đau không làm em vĩ đại, nhưng có thể làm em thật hơn.
\item Người từng vỡ thường nhận ra nỗi vỡ của người khác nhanh hơn.
\item Nếu chỉ ôm đau để tự thương hại thì nó sẽ thối đi.
\item Đau được viết lại thành bài học thì mới bớt độc.
\item Có những phần người chỉ nở ra sau khi đã từng mất một cái gì đó quan trọng.
\item Chiều sâu là khi em thôi xem mọi thứ chỉ theo kiểu được và mất ngay trước mắt.
\item Không ai cần phải cảm ơn mọi nỗi đau, nhưng có thể học cách không để nó phí.
\item Càng hiểu vết thương của mình càng bớt ném nó vào người khác.
\item Có những bài học phải rơi xuống mới nhặt được.
\item Sau mỗi lần đau tử tế hơn một chút đã là lời đáp đẹp với quá khứ.
\item Đừng khoe vết sẹo, hãy dùng nó để sống khôn và ấm hơn.
\item Đau không cho em quyền làm đau người khác.
\item Chiều sâu thật thường nói vừa đủ, không cần diễn cảnh từng trải suốt ngày.
\item Một người đi qua đủ chuyện mà vẫn giữ được lòng là người rất hiếm.
\item Chữa lành không xóa sẹo, nó dạy em không sống như sắp chảy máu lại liên tục.
\item Mỗi nỗi đau tiêu hóa được là một lớp vỏ cứng vừa đủ cho tâm hồn.
\item Đi qua nhiều chuyện rồi vẫn còn tin vào điều đúng là một dạng đẹp rất mạnh.
\item Đau không nên là cái cớ để sống cộc, nó nên là lý do để sống có chiều sâu hơn.
\end{enumerate}

\chapter{100 câu cục súc về tự cứu mình, sống sâu và đi đường dài}
\markboth{100 câu cục súc về tự cứu mình, sống sâu và đi đường dài}{100 câu cục súc về tự cứu mình, sống sâu và đi đường dài}

\section{25 câu về tự cứu mình trước tiên}
\begin{enumerate}[leftmargin=*]
\item Không ai cứu nổi em nếu em vẫn nghiện phá mình.
\item Người đầu tiên phải chịu trách nhiệm cho đời em vẫn là em.
\item Đừng đợi một biến cố đủ lớn mới chịu tỉnh.
\item Tự cứu không phải tự ôm hết, mà là thôi ngồi đó chờ phép màu.
\item Có những ngày tiến bộ là chỉ vì em không bỏ mặc mình như cũ.
\item Cần giúp đỡ là bình thường, nhưng phần bước chân đầu tiên vẫn phải do em nhấc.
\item Tự cứu bắt đầu từ những quyết định rất nhỏ nhưng không còn phản bội bản thân nữa.
\item Muốn người khác tin mình thay đổi thì trước hết mình phải thôi lặp lại điều cũ khi không ai thấy.
\item Đừng chỉ đọc thứ chữa lành, hãy sửa thứ đang làm em chảy máu.
\item Tự cứu là dọn lại cái đời sống thường ngày chứ không chỉ đổi vài câu nói với bản thân.
\item Có những người cầu cứu rất to nhưng chưa hề bỏ nổi điều đang đẩy họ xuống.
\item Tự cứu không cần hùng hồn, nó cần thật.
\item Hôm nào cũng đợi tâm trạng tốt mới sống tử tế với mình thì hơi lâu.
\item Người trưởng thành biết lúc nào cần nghỉ và lúc nào đang lười núp sau chữ nghỉ.
\item Tự cứu không đẹp, có khi nó chỉ là đi ngủ đúng giờ và ngừng nhắn tin cho người khiến mình nát.
\item Một quyết định đúng lặp lại đủ lâu sẽ cứu em mạnh hơn một cơn cảm hứng lớn.
\item Cứu mình không phải biến thành siêu nhân, chỉ là bớt làm kẻ phá hoại đời mình.
\item Em không cần hoàn hảo để bắt đầu hồi sinh.
\item Tự thương đúng là cho mình cơ hội sửa, không phải cho mình miễn hậu quả mãi.
\item Ai cũng muốn có điểm tựa, nhưng nhiều người quên mình cũng phải thành điểm tựa cho mình.
\item Điều khó nhất không phải biết mình cần làm gì, mà là làm nó khi chẳng ai vỗ tay.
\item Cứu mình là ngừng đặt chìa khóa bình yên vào tay người khác.
\item Đừng đợi một ngày đẹp trời, hãy dùng cái ngày đang có.
\item Mỗi lần em chọn điều bền hơn điều tiện là một lần tự cứu.
\item Cuộc đời khá lên từ lúc em ngừng giao quyền lái cho sự bốc đồng.
\end{enumerate}

\section{25 câu về sống sâu giữa thời đại ồn}
\begin{enumerate}[leftmargin=*,resume]
\item Thời đại này rất ồn, người có chiều sâu cần học cách im đúng lúc.
\item Đừng để mạng xã hội quyết định em phải buồn, vui, thành công và đẹp theo khuôn nào.
\item Nhiều thứ đang viral không có nghĩa nó đáng cho em đem đời mình ra học.
\item Càng nhiều nội dung nhanh càng cần vài khoảng chậm để não em không thành chợ.
\item Sống sâu không phải đăng câu gì cũng thấm, mà là thật sự nghĩ trước khi sống.
\item Người nông thường sợ sự im lặng vì không có gì nâng họ trong đó.
\item Đọc ít một mà ngấm còn hơn nuốt cả đống câu hay rồi sống y nguyên.
\item Đừng để thói quen lướt nhanh làm đầu óc em mất khả năng ở lại với điều quan trọng.
\item Nội lực được xây trong phần đời không ai quay clip.
\item Một người trẻ rất cần có chỗ trú khỏi tiếng ồn của mọi so sánh và kích thích liên tục.
\item Sống sâu là biết điều gì đáng giữ lâu hơn một cơn trend.
\item Không phải cái gì nhiều người xem cũng nhiều giá trị.
\item Tâm hồn đói chiều sâu rất dễ nghiện tiêu thụ cảm xúc rẻ.
\item Đừng nhầm việc biết nhiều khái niệm chữa lành với việc thật sự đang lành hơn.
\item Có khoảng lặng rồi em mới nghe được tiếng gì trong mình là thật.
\item Người biết dừng màn hình đúng lúc thường giữ được mắt nhìn đời rõ hơn.
\item Sống sâu nghĩa là không để bản thân bị giật qua lại bởi mọi thứ lấp lánh.
\item Đầu óc cần lọc rác như căn phòng cần dọn bụi.
\item Hãy giữ vài điều thiêng khỏi việc phải chứng minh với thiên hạ liên tục.
\item Chiều sâu cần thời gian nấu, không có bản instant đáng tin.
\item Càng nhiều tiếng ồn càng phải có lõi giá trị riêng.
\item Đừng để mình thành người chỉ biết phản ứng, hãy thành người biết chọn phản ứng nào đáng.
\item Một ngày không cần gì ghê gớm, chỉ cần đủ chậm để em còn nghe thấy mình.
\item Sống sâu là bớt dùng đời thật để nuôi hình ảnh ảo.
\item Người giữ được chiều sâu ở thời đại này là người biết tự bảo vệ sự tập trung như tài sản.
\end{enumerate}

\section{25 câu về đi đường dài}
\begin{enumerate}[leftmargin=*,resume]
\item Đi đường dài cần phổi, tim và cái đầu biết tính, không chỉ cần lửa.
\item Bùng lên dễ hơn duy trì nhiều.
\item Đừng ham thắng nhanh tới mức đốt sạch sức của tháng sau.
\item Những người đi xa ít khi sống toàn theo cơn hứng.
\item Đường dài rất ghét người thích cảm giác thành công hơn năng lực thật.
\item Cứ đều rồi một ngày em sẽ nhìn lại và thấy mình đã đi xa hơn kiểu ồn ào rất nhiều.
\item Bền không hào nhoáng, nhưng đánh bại được rất nhiều kiểu tài năng thất thường.
\item Đi đường dài nghĩa là biết lúc nào tăng, lúc nào nghỉ, lúc nào giữ nhịp.
\item Đừng khinh sự ổn định, nó là xương sống của nhiều thành tựu.
\item Người giỏi một đợt không đáng sợ bằng người biết giữ nhịp mười năm.
\item Mỗi lần em chọn làm đúng việc nhỏ là đang góp vào đời dài của mình.
\item Không phải ai khởi đầu mạnh cũng về đích đẹp.
\item Đời dài không dành cho người lúc nào cũng muốn thấy phần thưởng ngay.
\item Càng lớn càng phải biết quý năng lượng như quý tiền.
\item Đi lâu được là nhờ bớt tiêu hao vào mấy cuộc đọ vô nghĩa.
\item Đường dài đòi em bỏ vài thú vui ngắn nhưng đắt.
\item Thứ giữ em đi lâu không chỉ là mục tiêu, mà là cách sống đủ lành để không tự thiêu.
\item Người đi xa biết quý việc nhỏ lặp lại hơn mấy cơn lên đồng đầy cảm hứng.
\item Đừng đánh đổi sức khỏe tinh thần lấy vài thành quả nhìn choáng mà rất ngắn.
\item Đi đường dài cần lòng kiên nhẫn với chính phiên bản chưa hoàn chỉnh của mình.
\item Cứ làm điều đúng đủ lâu, hiệu ứng cộng dồn sẽ nói thay cho những ngày tưởng chẳng đi đâu.
\item Không phải hôm nào cũng thấy ý nghĩa, nhưng hướng đi đúng sẽ trả lời dần.
\item Đừng bỏ đường dài chỉ vì mấy tháng chưa sáng rỡ.
\item Người bền không luôn chiến thắng hôm nay, nhưng họ hiếm khi tự loại mình khỏi cuộc chơi.
\item Đời không chỉ cần em cháy, nó cần em sáng lâu.
\end{enumerate}

\section{25 câu về một đời đáng sống}
\begin{enumerate}[leftmargin=*,resume]
\item Một đời đáng sống không phải đời dễ, mà là đời em hiểu mình đang sống vì gì.
\item Đừng sống chỉ để né đau, em sẽ né luôn nhiều điều đáng giá.
\item Giá trị của đời em không đo hết bằng mức độ người khác công nhận.
\item Sống cho ra người đôi khi bắt đầu từ việc bớt làm điều khiến mình xấu đi.
\item Đừng đợi cuộc đời trả lời hết mới chịu bước.
\item Có ý nghĩa không phải lúc nào cũng vui, nhưng thường làm em đỡ rỗng.
\item Một đời đẹp không cần ồn nếu bên trong có trật tự và tình người.
\item Đừng để nỗi sợ bình thường làm em chọn sống nhạt với mọi thứ.
\item Điều đáng tiếc nhất không phải là chưa nổi tiếng, mà là sống mà chẳng thành mình tử tế hơn.
\item Có tiền, có việc, có người thương đều quý, nhưng có cốt cách để giữ những thứ đó còn quý hơn.
\item Sống kỹ một chút để đến lúc nhìn lại đỡ thấy mình toàn chạy theo bụi.
\item Một đời đáng sống là đời em không cần liên tục trốn khỏi chính mình.
\item Bình yên thật không nằm ở việc chẳng có sóng, mà là em có tay lái.
\item Đừng để mình chỉ tồn tại, hãy tập sống có lựa chọn.
\item Càng lớn càng hiểu điều quý nhất không phải lúc nào cũng là điều ồn nhất.
\item Đừng mua cả đời bằng vài khoảnh khắc được người khác trầm trồ.
\item Sống đẹp không làm em thành thánh, nó chỉ giúp em ngủ yên hơn.
\item Hãy làm người mà khi mệt nhất vẫn không thấy ghê với chính mình.
\item Nếu một ngày không còn ai nhìn, em vẫn muốn thành người như thế nào?
\item Trưởng thành là biết cái gì đáng mang theo lâu, cái gì nên bỏ bên đường.
\item Đời đáng sống không phải đời không sai, mà là đời biết sửa sai mà không mất tim.
\item Cuối cùng, người cứu em, giữ em và làm đời em ra gì vẫn sẽ là những lựa chọn lặp lại của em.
\item Một người sâu không cần nhiều khẩu hiệu, chỉ cần sống bớt giả hơn hôm qua.
\item Nếu phải chọn, hãy chọn thành người có ích mà vẫn còn phần người.
\item Sống sao để khi quay lại, em thấy mình đã đi qua đời này không quá uổng.
\end{enumerate}